\documentclass[journal,10pt,twocolumn]{IEEEtran}

\usepackage{amsmath,amsfonts,amssymb}
\usepackage{graphicx}
\usepackage{booktabs}
\usepackage{cite}
\usepackage{url}
\usepackage{microtype}
\usepackage{tabularx}
\usepackage{multirow}
\usepackage{xcolor}

\begin{document}

\title{AstraVisionX: A Multimodal Vision Platform with Multi-Method Explainability, Consensus-Based Reliability Analysis, and Visually Grounded Conversational Reasoning}

\author{Gaurav~Kumar$^1$,
        Shail~Rawat$^2$
\thanks{$^{1,2}$School of Engineering and Technology (SOET), K.R. Mangalam University, Gurugram, Haryana, India. (e-mail: gkc1764@gmail.com, shailrawat70@gmail.com).}}

\markboth{Preprint: Multi-Method Vision and XAI Evaluation}{Kumar and Rawat: AstraVisionX Multimodal Vision and XAI Platform}

\maketitle

\begin{abstract}
Deep computer vision architectures have achieved remarkable performance on object localization and scene classification benchmarks. However, deployed systems frequently operate as black boxes, yielding predictions without indicating why a decision was reached, how robust the underlying decision boundary is, or whether distinct explainable AI (XAI) algorithms agree. This paper presents \textbf{AstraVisionX}, an end-to-end multimodal computer vision platform integrating: (1) fine-grained object detection via YOLOv8 trained on Open Images V7; (2) global scene classification via a Vision Transformer (ViT-B/16); (3) a triple-method XAI suite comprising Grad-CAM, Integrated Gradients, and LIME; (4) a quantitative explanation consensus engine evaluating pairwise and global Intersection over Union (IoU) and Dice similarity across salient binary masks against analytical chance and empirical spatial nulls; (5) decision-boundary margin reliability estimation ($m = p_1 - p_2$) evaluated via Expected Calibration Error (ECE) and AUROC; and (6) a visually grounded multimodal conversational assistant powered by a multi-model cascade failover mechanism. \textit{AstraVisionX is intended for research and educational evaluation and is not a clinical diagnostic tool.}

Through an empirical evaluation across a 100-class ImageNet-1K benchmark, we observe a global explanation agreement of $A_{\text{global}} = 21.33\% \pm 1.44\%$ IoU ($33.97\% \pm 1.90\%$ Dice) at the top-20\% salient area threshold. While significantly exceeding the uniform random chance baseline ($11.11\%$), this overlap falls slightly below the empirical cross-image spatial null ($26.30\%$) induced by natural photographic center bias. For equal-area masks, this indicates that approximately $66\%$ of salient pixels identified by one explainer are unshared by another ($1 - \text{Dice}$), empirically substantiating the ``Disagreement Problem in XAI'' and demonstrating that single-method interpretability is fundamentally precarious. Faithfulness evaluations via deletion and insertion curves reveal that Integrated Gradients attains a deletion AUC of $0.152$ compared to $0.204$ for random masking, while Grad-CAM achieves high insertion sufficiency ($0.701$). Furthermore, cascading parameter randomization demonstrates that attributions collapse as randomization proceeds from classification heads ($\rho = 0.304$) into convolutional blocks ($\rho = -0.022$). The decision margin $m = p_1 - p_2$ achieves an AUROC of $0.998$ ($95\%$ bootstrap CI: $[0.985, 1.000]$) for misclassification detection, with accuracy dropping from $100.0\%$ for $m \ge 0.50$ to $16.7\%$ for $m < 0.20$. The platform is implemented as an asynchronous FastAPI microservice with a responsive glassmorphic dashboard, client-side Web Speech interaction, and automated PDF dossier generation.
\end{abstract}

\begin{IEEEkeywords}
Explainable AI, Grad-CAM, Integrated Gradients, LIME, Explanation Consensus, Disagreement Problem, Faithfulness Evaluation, Decision Boundary, Vision Transformer, Multimodal LLM.
\end{IEEEkeywords}

\section{Introduction}
\IEEEPARstart{C}{omputer} vision models based on deep convolutional networks and self-attention transformers have attained impressive accuracy across diverse visual benchmarks. Despite these successes, deep neural networks remain opaque function approximators. In high-stakes applications—including autonomous navigation, industrial quality inspection, and scientific analysis—a point prediction in isolation is insufficient~\cite{selvaraju2017gradcam}. Domain practitioners must comprehend which visual features influenced a prediction, whether the model's decision boundary is well-separated, and whether disparate interpretability techniques corroborate one another. \textit{AstraVisionX is designed as a research and educational platform and is explicitly not a clinical diagnostic system.}

Explainable Artificial Intelligence (XAI) has emerged to address model opacity by generating post-hoc attribution heatmaps. Popular paradigms include gradient-based activation mapping (e.g., Grad-CAM~\cite{selvaraju2017gradcam}), axiomatic path integration (e.g., Integrated Gradients~\cite{sundararajan2017axiomatic}), and perturbation-based local surrogates (e.g., LIME~\cite{ribeiro2016why}). However, real-world deployment of XAI faces three foundational challenges:
\begin{enumerate}
    \item \textbf{The Disagreement Problem}: Distinct XAI algorithms frequently attribute high importance to non-overlapping image regions for the identical input and model~\cite{krishna2022disagreement}. Relying on a single explainer exposes practitioners to confirmation bias.
    \item \textbf{Uncalibrated Decision Boundaries}: Softmax output probabilities are prone to calibration gaps on ambiguous inputs~\cite{guo2017calibration,hendrycks2017baseline}. High raw confidence does not imply decision stability.
    \item \textbf{Static, Non-Interactive Visualizations}: Traditional XAI systems generate passive 2D heatmap plots that cannot be queried, segmented per detected object, or interactively reasoned over by downstream users.
\end{enumerate}

To address these challenges, we propose \textbf{AstraVisionX}, an open-source multimodal vision and reasoning platform. Given an arbitrary visual input, AstraVisionX executes fine-grained object detection, scene classification, triple-method visual explanation, quantitative consensus analysis, decision-boundary reliability estimation, and conversational multimodal reasoning.

\noindent\textbf{Contributions.} The primary contributions of this work are:
\begin{enumerate}
    \item \textbf{Unified Multi-Method XAI Framework}: We integrate Grad-CAM, Integrated Gradients, and LIME over a shared, controlled backbone to systematically evaluate attribution convergence across algorithmic paradigms.
    \item \textbf{Quantitative Consensus Metric with Chance and Spatial Nulls}: We formulate an automated consensus engine that binarizes continuous heatmaps using top-$k\%$ salient area thresholding and computes pairwise/global IoU and Dice similarity against closed-form chance baselines ($\mathbb{E}[\text{IoU}_{\text{chance}}] = \frac{k}{2-k}$) and empirical cross-image spatial nulls.
    \item \textbf{Empirical Faithfulness and Reliability Validation}: We evaluate attribution faithfulness via deletion/insertion curves~\cite{petsiuk2018rise} and cascading parameter randomization~\cite{adebayo2018sanity}, while validating decision margin reliability ($m = p_1 - p_2$) via Expected Calibration Error (ECE) and AUROC for misclassification detection.
    \item \textbf{Visually Grounded Multimodal Reasoning}: We ground a multimodal conversational agent with structured bounding boxes, classification distributions, and attribution metrics, describing a standardized visual question-answering workflow.
    \item \textbf{Scalable Full-Stack Architecture}: We provide an asynchronous FastAPI backend featuring an in-memory TTL session manager, a multi-model failover engine, and an interactive browser dashboard with Web Speech support.
\end{enumerate}

\section{Related Work}
\subsection{Object Detection and Scene Classification}
Modern object detection frameworks balance inference speed and localization fidelity. The YOLO family~\cite{jocher2023yolo} provides anchor-free single-stage detection. Trained on the Open Images V7 dataset~\cite{kuznetsova2020openimages}, YOLOv8 models localize over 600 common object categories with high mean Average Precision. In parallel, Vision Transformers (ViT)~\cite{dosovitskiy2021vit} divide visual inputs into non-overlapping spatial patches, processing long-range spatial dependencies through self-attention layers to achieve state-of-the-art ImageNet-1K classification~\cite{deng2009imagenet}.

\subsection{Visual Explanation Methods and Disagreement}
Post-hoc visual explainability encompasses gradient-weighted feature map pooling (Grad-CAM~\cite{selvaraju2017gradcam}), path-integrated gradients satisfying implementation invariance and completeness (Integrated Gradients~\cite{sundararajan2017axiomatic}), and interpretable local surrogate models (LIME~\cite{ribeiro2016why}). SHAP~\cite{lundberg2017unified} extends Shapley value attribution to neural networks. 

Crucially, recent studies have highlighted the \textit{Disagreement Problem in XAI}~\cite{krishna2022disagreement}: explainers operating on different mathematical premises often exhibit modest spatial consensus. While software toolkits such as Captum~\cite{kokhlikyan2020captum}, Quantus~\cite{hedstrom2023quantus}, and OmniXAI~\cite{yang2022omnixai} provide developer libraries for computing individual attributions, to our knowledge they lack an integrated, real-time consensus engine that exposes agreement metrics to end users in an interactive dashboard.

\subsection{Multimodal Conversational Agents}
Vision-language foundation models (e.g., LLaVA~\cite{liu2023visual}, Flamingo~\cite{alayrac2022flamingo}, and Gemini~\cite{geminiteam2024gemini}) enable natural language dialogue grounded in pixel inputs. However, vanilla vision-language models remain prone to spatial approximations. AstraVisionX addresses this vulnerability by injecting structured detection coordinates, classification margins, and XAI consensus metrics directly into the multimodal context window.

\section{Proposed AstraVisionX Framework}

\subsection{System Architecture}
AstraVisionX is structured into five cohesive architectural layers:
\begin{enumerate}
    \item \textbf{Presentation Layer}: Browser dashboard supporting drag-and-drop ingestion, pan/zoom viewports, in-place XAI crossfading, and Web Speech voice interaction.
    \item \textbf{Orchestration Layer}: Asynchronous FastAPI microservice managing thread-safe sessions with 30-minute Time-To-Live (TTL) cache eviction.
    \item \textbf{Inference Layer}: Dual-engine execution comprising Ultralytics YOLOv8s object detection and timm ViT-B/16 classification.
    \item \textbf{Explanation and Analysis Layer}: Grad-CAM, Integrated Gradients, and LIME generators evaluated by a pairwise IoU/Dice consensus calculator.
    \item \textbf{Reasoning and Reporting Layer}: Multi-key, multi-model failover Gemini conversational assistant and ReportLab PDF dossier compiler.
\end{enumerate}

\subsection{Dual Deep Learning Pipeline}
\noindent\textbf{Fine-Grained Object Detection and Hierarchical Filtering.} We deploy YOLOv8s trained on Open Images V7 (601 classes). Given an RGB input image $I \in \mathbb{R}^{H \times W \times 3}$, the detector infers $N$ bounding boxes:
\begin{equation}
B_n = [x_{\min}, y_{\min}, x_{\max}, y_{\max}], \quad c_n \in [0, 1], \quad y_n \in \mathcal{C}_{\text{OI}}
\end{equation}
where $c_n$ represents confidence and $\mathcal{C}_{\text{OI}}$ denotes the 601 Open Images classes. To eliminate channel-swapping distortions, arrays are explicitly calibrated from BGR to RGB prior to tensor transformation. Furthermore, detections are systematically mapped into top-level semantic categories ($\mathcal{T} = \{\text{Persons}, \text{Vehicles}, \text{License Plates}, \text{Clothing}, \text{Animals}, \text{Other}\}$), with parent-child hierarchical containment linking sub-detections (e.g., worn clothing to person; registration plate to vehicle) based on bounding-box containment intersection-over-area ($\text{IoA} \ge 0.50$).

\noindent\textbf{Global Scene Classification.} Simultaneously, global visual semantics are extracted via a Vision Transformer ($\text{ViT-Base-Patch16-224}$) pretrained on ImageNet-1K ($\mathcal{C}_{\text{IN}}$, 1000 classes). Input images are resized to $224 \times 224$ and projected through patch self-attention blocks to produce class probabilities $\mathbf{p} = \text{Softmax}(\mathbf{z})$.

\noindent\textbf{Reconciliation of Label Spaces.} Because $\mathcal{C}_{\text{OI}}$ and $\mathcal{C}_{\text{IN}}$ are distinct label taxonomies, AstraVisionX decouples regional detection from global classification. Furthermore, to enable localized interpretability, AstraVisionX supports \textit{Region-Targeted XAI}: the user can select any detected object $B_n$, which is dynamically cropped with $20\%$ contextual boundary padding, explained via the XAI suite, and smoothly blended back onto the full canvas using a Gaussian boundary smoothing kernel ($k = \max(5, 0.15 \min(w, h))$).

\subsection{Triple-Method Explainable AI Suite}
To guarantee a strictly controlled attribution benchmark, all three XAI explainers operate on a shared ResNet-18~\cite{he2016deep} backbone $F(x)$, ensuring that mathematical differences in explanation maps reflect algorithmic attribution properties rather than architectural disparities.

\subsubsection{Grad-CAM}
Grad-CAM computes class-discriminative localization by pooling gradients flowing into the final convolutional layer (\texttt{layer4.1.conv2}). The neuron importance weight $\alpha_k^c$ for feature map $A^k$ with respect to class score $y^c$ is:
\begin{equation}
\alpha_k^c = \frac{1}{Z} \sum_{i=1}^U \sum_{j=1}^V \frac{\partial y^c}{\partial A_{ij}^k}
\end{equation}
The heat activation map is obtained via a rectified linear combination:
\begin{equation}
L_{\text{Grad-CAM}}^c = \text{ReLU}\left(\sum_k \alpha_k^c A^k\right)
\end{equation}
which is upsampled to input resolution $(H, W)$ using bilinear interpolation.

\subsubsection{Integrated Gradients}
Integrated Gradients computes path-integrals along straight-line trajectories between a baseline reference $x'$ and input $x$:
\begin{equation}
\label{eq:ig_continuous}
IG_i(x) = (x_i - x_i') \times \int_0^1 \frac{\partial F(x' + \alpha(x - x'))}{\partial x_i} d\alpha
\end{equation}
We approximate \eqref{eq:ig_continuous} via a Riemann summation across $m = 8$ interpolation steps. To mitigate baseline selection artifacts, AstraVisionX evaluates multiple baselines (black zero and dataset empirical mean), aggregating attributions across references:
\begin{multline}
IG_i^{\text{multi}}(x) = \frac{1}{|B|} \sum_{x' \in B} (x_i - x_i') \\
\times \frac{1}{m} \sum_{k=1}^m \frac{\partial F(x' + \frac{k}{m}(x - x'))}{\partial x_i}
\end{multline}

\subsubsection{LIME}
LIME segments the image into $K = 20$ superpixels using Simple Linear Iterative Clustering (SLIC)~\cite{achanta2012slic} with compactness $10$ and Gaussian smoothing $\sigma = 1$. A binary perturbation matrix $Z \in \{0, 1\}^{N \times K}$ ($N=25$) is generated, where $z_k = 0$ replaces superpixel $k$ with mean pixel values. A ridge regression surrogate model $g(z') = w^T z'$ is minimized with an exponential distance kernel $\pi(z) = \exp(-D(z, \mathbf{1})^2 / \sigma_{\text{dist}}^2)$ ($\sigma_{\text{dist}} = 0.25$).

\subsection{Explanation Consensus Analysis}
To provide a scale-invariant comparison across continuous heatmaps $M \in \mathbb{R}^{H \times W}$, we define the salient mask $\hat{M}_k$ using a \textbf{Top-$k\%$ salient area threshold}:
\begin{equation}
\tau_k = \text{Percentile}_{100-k}(M)
\end{equation}
\begin{equation}
\hat{M}_{k, ij} = \mathbb{I}(M_{ij} \ge \tau_k)
\end{equation}
This formulation targets an identical salient area fraction $k \in (0, 1)$ across explainers (subject to minor tie effects and discrete superpixel boundaries). For two binary masks $A$ and $B$, the Intersection over Union (IoU) and Dice similarity are:
\begin{equation}
\text{IoU}(A, B) = \frac{|A \cap B|}{|A \cup B|}, \quad \text{Dice}(A, B) = \frac{2|A \cap B|}{|A| + |B|}
\end{equation}
When two independent random masks each select a fraction $k$ of pixels, the expected random overlap yields a closed-form \textbf{Chance Baseline}:
\begin{equation}
\mathbb{E}[\text{IoU}_{\text{chance}}] = \frac{k^2}{2k - k^2} = \frac{k}{2 - k}, \quad \mathbb{E}[\text{Dice}_{\text{chance}}] = k
\end{equation}
Across Grad-CAM ($G$), Integrated Gradients ($I$), and LIME ($L$), the Global Agreement Score $A_{\text{global}}$ is:
\begin{equation}
A_{\text{global}} = \frac{\text{IoU}(G, I) + \text{IoU}(G, L) + \text{IoU}(I, L)}{3}
\end{equation}
For equal-cardinality masks, the proportion of one mask's pixels unshared by another is exactly $1 - \text{Dice}(A, B)$, whereas $1 - \text{IoU}(A, B)$ represents the Jaccard distance.

\subsection{Decision-Boundary Reliability Analysis}
Given top predicted probability $p_1$ and runner-up probability $p_2$, we compute the decision-boundary separation margin:
\begin{equation}
m = p_1 - p_2
\end{equation}
AstraVisionX categorizes predictions into three operational tiers:
\begin{itemize}
    \item \textbf{High Reliability}: $m \ge 0.50$ (decisive boundary separation).
    \item \textbf{Moderate Reliability}: $0.20 \le m < 0.50$.
    \item \textbf{Narrow Decision Boundary}: $m < 0.20$ (potential boundary ambiguity).
\end{itemize}
Detections below threshold are explicitly relabeled with warning badges (e.g., ``$\triangle$ Uncertain match'').

\subsection{Multimodal Conversational Reasoning}
AstraVisionX integrates Google Gemini foundation models~\cite{geminiteam2024gemini} via a multi-key load balancer and multi-model cascade failover client (\texttt{gemini\_service.py}). When rate-limit (HTTP 429) or high-demand (HTTP 503) errors arise, requests automatically cascade across prioritized candidate models (\texttt{gemini-3.1-flash-lite}, \texttt{gemini-3.7-flash}, \texttt{gemini-3.5-flash}), caching the active key/model for sub-second subsequent invocations. 

The assistant receives a dual input: (1) an optimized 512px thumbnail JPEG, and (2) structured session text containing bounding box coordinates, top-5 classification rankings, XAI agreement percentages, and boundary margins.

\section{System Implementation}
The backend is built in Python 3.14 using FastAPI and Uvicorn. All state is maintained in an in-memory thread-safe \texttt{SessionManager} protected by re-entrant mutexes with 30-minute automated expiration. The REST interface exposes \texttt{/health}, \texttt{/detect}, \texttt{/explain/method}, \texttt{/analyze}, \texttt{/report}, and \texttt{/chat}. The frontend comprises vanilla HTML5, CSS3 glassmorphic design tokens, ES6+ JavaScript, and the browser-native Web Speech API. Comprehensive analysis dossiers are compiled programmatically into downloadable PDFs using ReportLab.

\section{Experimental Protocol}
\noindent\textbf{Evaluation Benchmark Dataset.} Experiments were conducted over a 100-class evaluation benchmark sampled uniformly (every tenth category) across the full 1,000-class ImageNet-1K validation set, alongside Open Images V7 object detection benchmarks.

\noindent\textbf{Hardware Environment.} All benchmarks were executed on an Intel/AMD workstation CPU (single-precision FP32 PyTorch 2.x execution, 16\,GB RAM, Windows 11 x86\_64) to establish a baseline for unaccelerated edge and cloud deployments.

\section{Results and Empirical Analysis}

\subsection{Detection and Classification Baselines (Table~\ref{tab:baselines})}
Table~\ref{tab:baselines} summarizes baseline benchmark performance. YOLOv8s achieves $0.448$ mAP@0.5 across 601 Open Images classes, while ViT-B/16 achieves $84.2\%$ published Top-1 accuracy on ImageNet-1K. On our curated 100-class benchmark of clear category exemplars, ViT-B/16 achieves an empirical Top-1 accuracy of $94.0\%$.

\begin{table}[t]
\centering
\caption{Detection and classification baseline performance.}
\label{tab:baselines}
\begin{tabular}{lcccc}
\toprule
\textbf{Model} & \textbf{mAP@.5} & \textbf{mAP@.5:.95} & \textbf{Top-1} & \textbf{Top-5} \\
\midrule
YOLOv8s (Open Images V7) & 0.448 & 0.312 & -- & -- \\
ViT-B/16 (ImageNet-1K published) & -- & -- & 84.2\% & 97.1\% \\
ViT-B/16 (100-class evaluated) & -- & -- & 94.0\% & 99.0\% \\
ResNet-18 (XAI Backbone) & -- & -- & 69.8\% & 89.1\% \\
\bottomrule
\end{tabular}
\end{table}

\subsection{Explanation Agreement and the Disagreement Problem (Table~\ref{tab:agreement})}
Table~\ref{tab:agreement} reports pairwise and global agreement scores across the $N=100$ benchmark dataset at the top-20\% salient area threshold ($k=20\%$), accompanied by $95\%$ confidence intervals ($\text{CI}_{95}$) and the analytical chance baseline.

\begin{table}[t]
\centering
\caption{Quantitative multi-method explanation agreement ($N=100$, top-$20\%$ salient area threshold, with $95\%$ confidence intervals and chance baselines).}
\label{tab:agreement}
\begin{tabular}{lcccc}
\toprule
\textbf{Method Pair} & \textbf{Mean IoU (\%)} & \textbf{Median} & \textbf{Chance} & \textbf{Mean Dice (\%)} \\
\midrule
Grad-CAM -- IG & $23.40 \pm 2.21$ & 21.92\% & 11.11\% & $36.63 \pm 2.81$ \\
Grad-CAM -- LIME & $23.05 \pm 2.32$ & 22.50\% & 11.11\% & $36.01 \pm 3.02$ \\
IG -- LIME & $17.55 \pm 1.36$ & 16.90\% & 11.11\% & $29.28 \pm 1.94$ \\
\midrule
\textbf{Global ($A_{\text{global}}$)} & $\mathbf{21.33 \pm 1.44}$ & \textbf{20.43\%} & \textbf{11.11\%} & $\mathbf{33.97 \pm 1.90}$ \\
\bottomrule
\end{tabular}
\end{table}

While the empirical global agreement ($21.33\% \pm 1.44\%$) exceeds the uniform random chance baseline ($11.11\%$), an evaluation of empirical cross-image pairings reveals a spatial null overlap of $26.30\%$ driven by natural photographic center bias. Because natural subject matter concentrates near image centers, two arbitrary images of distinct classes exhibit $\approx 26.3\%$ mask overlap. Hence, cross-method consensus ($21.33\%$) is comparable to or lower than the empirical spatial center-bias null. For equal-cardinality masks, a mean Dice of $33.97\%$ indicates that approximately $66\%$ of salient pixels identified by one explainer are unshared by another, empirically substantiating the Disagreement Problem in XAI~\cite{krishna2022disagreement}.

\noindent\textbf{Threshold Sensitivity Analysis.} Table~\ref{tab:sensitivity} reports agreement across salient area thresholds $k \in \{10\%, 15\%, 20\%, 25\%, 30\%\}$. Across all cut-offs, empirical IoU tracks approximately $9.5\% - 11\%$ above the mathematical uniform chance baseline, demonstrating that explanation divergence is structural rather than an artifact of threshold selection.

\begin{table}[t]
\centering
\caption{Threshold sensitivity analysis across top-$k\%$ salient area ($N=100$).}
\label{tab:sensitivity}
\begin{tabular}{lcccc}
\toprule
\textbf{Salient Area ($k$)} & \textbf{Global IoU (\%)} & \textbf{$\text{CI}_{95}$} & \textbf{Chance IoU} & \textbf{GC vs. IG (\%)} \\
\midrule
Top 10\% & 14.73\% & $\pm 1.37$ & 5.26\% & 15.59\% \\
Top 15\% & 17.78\% & $\pm 1.39$ & 8.11\% & 19.77\% \\
\textbf{Top 20\% (Default)} & \textbf{21.33\%} & $\mathbf{\pm 1.44}$ & \textbf{11.11\%} & \textbf{23.40\%} \\
Top 25\% & 24.87\% & $\pm 1.34$ & 14.29\% & 27.37\% \\
Top 30\% & 28.76\% & $\pm 1.49$ & 17.65\% & 30.37\% \\
\bottomrule
\end{tabular}
\end{table}

\subsection{Explanation Faithfulness: Deletion and Insertion Curves (Table~\ref{tab:faithfulness})}
Following the methodology of Petsiuk et al.~\cite{petsiuk2018rise}, we measured explanation faithfulness and sufficiency via deletion and insertion area-under-the-curve (AUC) metrics across progressive perturbation steps (mean-pixel baseline, $4$ steps, ResNet-18 top-1 class). A faithful explanation exhibits a low Deletion AUC (rapidly degrading class confidence when salient pixels are masked) and a high Insertion AUC. 

\begin{table}[t]
\centering
\caption{Explanation faithfulness: Deletion and Insertion AUC~\cite{petsiuk2018rise}.}
\label{tab:faithfulness}
\begin{tabular}{lcc}
\toprule
\textbf{Attribution Method} & \textbf{Deletion AUC ($\downarrow$)} & \textbf{Insertion AUC ($\uparrow$)} \\
\midrule
Random Baseline & 0.204 & 0.214 \\
Grad-CAM & 0.296 & \textbf{0.701} \\
Integrated Gradients & \textbf{0.152} & 0.545 \\
\bottomrule
\end{tabular}
\end{table}

As reported in Table~\ref{tab:faithfulness}, Integrated Gradients attains a low Deletion AUC of $0.152$ (outperforming random attribution at $0.204$), indicating that removing its top-attributed pixels rapidly degrades target probability. Conversely, Grad-CAM exhibits a higher Deletion AUC ($0.296$) than random baseline, reflecting its coarse spatial activation nature: removing upsampled feature blobs can leave contextual object boundaries intact during initial steps. However, Grad-CAM achieves the highest insertion sufficiency ($0.701$), rapidly restoring prediction confidence when its salient region is introduced onto a blank canvas.

\subsection{Decision-Boundary Reliability and Calibration (Table~\ref{tab:reliability})}
On our evaluation suite, the Vision Transformer exhibits an Expected Calibration Error of $\text{ECE} = 0.1586$. With an empirical accuracy of $94.0\%$ and a mean confidence of $\approx 0.83$, the model displays slight underconfidence on these clear exemplars, unlike out-of-distribution inputs where overconfidence typically dominates.

When evaluated as a misclassification detector (positive class = correctly classified sample), the decision margin $m = p_1 - p_2$ achieves an $\text{AUROC} = 0.998$ ($95\%$ bootstrap CI: $[0.985, 1.000]$) and an $\text{AUPR} = 1.000$. As detailed in Table~\ref{tab:reliability}, samples in the High Margin tier ($m \ge 0.50$, mean margin $0.838$) and Moderate Margin tier ($0.20 \le m < 0.50$, mean margin $0.362$) both attain $100.0\%$ Top-1 accuracy. In contrast, samples in the narrow boundary tier ($m < 0.20$, mean margin $0.118$) drop sharply to $16.7\%$ accuracy. This empirical divergence provides operational justification for the $m = 0.20$ separation boundary.

\begin{table}[t]
\centering
\caption{Validation of decision-boundary margin $m$ vs. accuracy ($N=100$).}
\label{tab:reliability}
\begin{tabular}{lccc}
\toprule
\textbf{Reliability Tier} & \textbf{Count} & \textbf{Top-1 Accuracy} & \textbf{Mean Margin ($m$)} \\
\midrule
High Margin ($m \ge 0.50$) & 85 & 100.0\% & 0.838 \\
Moderate Margin ($0.20 \le m < 0.50$) & 9 & 100.0\% & 0.362 \\
Narrow Boundary ($m < 0.20$) & 6 & 16.7\% & 0.118 \\
\bottomrule
\end{tabular}
\end{table}

\subsection{Processing Latency (Table~\ref{tab:latency})}
Table~\ref{tab:latency} details execution latency across $N=100$ evaluations on an unaccelerated CPU workstation. Object detection and scene classification execute in under $0.25$s each. Grad-CAM executes in $0.135$s, while path-integrated and superpixel explainers require $1.37\text{s} - 1.46\text{s}$. The deterministic local analytical pipeline completes in $3.480\text{s}$. Remote multimodal reasoning via Gemini involves external API latency ($\approx 1.2\text{s} - 2.8\text{s}$) and is managed via the failover cascade client.

\begin{table}[t]
\centering
\caption{Processing latency across pipeline stages on CPU ($N=100$).}
\label{tab:latency}
\begin{tabular}{lc}
\toprule
\textbf{Pipeline Stage} & \textbf{Mean $\pm$ Std Latency (s)} \\
\midrule
YOLOv8s Object Detection & 0.241 $\pm$ 0.025 \\
ViT-B/16 Scene Classification & 0.240 $\pm$ 0.027 \\
Grad-CAM Activation Mapping & 0.135 $\pm$ 0.023 \\
Integrated Gradients ($m=8$) & 1.458 $\pm$ 0.148 \\
LIME (SLIC Superpixels, $N=25$) & 1.372 $\pm$ 0.323 \\
Consensus \& Agreement Metrics & 0.001 $\pm$ 0.001 \\
PDF Report Dossier Generation & 0.033 $\pm$ 0.029 \\
\midrule
\textbf{Total Local Pipeline Runtime} & \textbf{3.480 s} \\
\bottomrule
\end{tabular}
\end{table}

\subsection{Adebayo et al. Cascading Parameter Randomization Test}
Following the sanity check methodology of Adebayo et al.~\cite{adebayo2018sanity}, we performed cascading parameter randomization of the ResNet-18 backbone. Randomizing only the final linear classification layer (\texttt{fc}) caused Grad-CAM rank correlation with the trained map to drop to Spearman $\rho = 0.304$ ($\text{SSIM} = 0.463$, mask $\text{IoU} = 26.3\%$). When cascading randomization reached the final convolutional block (\texttt{layer4}), correlation collapsed to Spearman $\rho = -0.022$ and mask overlap fell to $12.5\%$ (comparable to the $11.11\%$ analytical chance baseline). Similarly, Integrated Gradients dropped to $\rho = 0.501$ upon \texttt{fc} randomization. This confirms that the explainers pass the cascading randomization sanity test and reflect learned representations.

\subsection{Grounded Assistant Qualitative Demonstration}
In conversational queries concerning object localization and feature importance, structured grounding (injecting bounding box coordinates and XAI statistics into the Gemini context) allows the assistant to quote exact spatial boundaries ($[x_{\min}, y_{\min}, x_{\max}, y_{\max}]$) and cite specific attribution percentages. In contrast, ungrounded prompts default to qualitative approximations. Rigorous human-subject evaluation across blind expert raters remains planned future work.

\section{Discussion}
The experimental results yield three critical insights. First, the modest spatial overlap observed across Grad-CAM, Integrated Gradients, and LIME ($21.33\%$), falling below the empirical center-bias spatial null ($26.30\%$), indicates that visual explanations cannot be treated as monolithic ground truths; evaluating multi-method consensus allows practitioners to detect when attribution maps are method-dependent artifacts. Second, decision-boundary margins $m = p_1 - p_2$ provide an immediate indicator of classification stability ($\text{AUROC} = 0.998$), successfully identifying ambiguous inputs where accuracy drops to $16.7\%$. Third, injecting structured coordinates and XAI statistics into vision-language foundation models provides grounding that helps mitigate visual hallucinations during dialogue.

\section{Limitations and Future Work}
Several limitations should be noted. First, the evaluation in this study was conducted over a 100-class benchmark sample; evaluating full 50,000-image datasets remains an objective for distributed compute clusters. Second, all XAI explanations in the current implementation describe the ResNet-18 convolutional backbone rather than the Vision Transformer or YOLOv8. Third, high explanation agreement does not guarantee attribution faithfulness, as demonstrated by Grad-CAM's higher deletion AUC relative to random masking. Fourth, LIME regressions with small sample sizes ($N=25$) exhibit sensitivity to perturbation noise. Finally, conversational reasoning depends on external Google Gemini APIs, transmitting image thumbnails externally. Future work will investigate attention-rollout explainability for Vision Transformers, local open-weight visual language models (e.g., LLaVA), and fine-tuning on specialized medical datasets.

\section{Conclusion}
This paper introduced AstraVisionX, an asynchronous multimodal vision platform integrating YOLOv8 object detection, Vision Transformer scene classification, triple-method explainability (Grad-CAM, Integrated Gradients, LIME), quantitative IoU/Dice consensus metrics against chance and spatial nulls, decision margin reliability estimation, and grounded conversational reasoning. By demonstrating that cross-method explanation agreement is inherently modest ($A_{\text{global}} = 21.33\%$) and that classification margins directly predict correctness ($\text{AUROC} = 0.998$), AstraVisionX transforms black-box computer vision into an interactive, reliable, and auditable visual analysis workflow.

\section*{Code and Data Availability}
The full AstraVisionX platform source code, pre-trained weights, benchmark evaluation dataset, and reproduction scripts are openly available at: \url{https://github.com/GauravKumar/AstraVisionX}.

\begin{thebibliography}{10}

\bibitem{selvaraju2017gradcam}
R.~R. Selvaraju, M.~Cogswell, A.~Das, R.~Vedantam, D.~Parikh, and D.~Batra,
  ``Grad-cam: Visual explanations from deep networks via gradient-based
  localization,'' in \emph{Proc. IEEE International Conference on Computer Vision (ICCV)}, 2017, pp. 618--626.

\bibitem{sundararajan2017axiomatic}
M.~Sundararajan, A.~Taly, and Q.~Yan, ``Axiomatic attribution for deep
  networks,'' in \emph{Proc. International Conference on Machine Learning (ICML)}, 2017, pp. 3319--3328.

\bibitem{ribeiro2016why}
M.~T. Ribeiro, S.~Singh, and C.~Guestrin, ``"Why should I trust you?":
  Explaining the predictions of any classifier,'' in \emph{Proc. ACM SIGKDD International Conference on Knowledge Discovery and Data Mining}, 2016, pp. 1135--1144.

\bibitem{krishna2022disagreement}
S.~Krishna, T.~Han, A.~Gu, J.~Poursabzi-Sangdeh, S.~Jabbari, P.~Wu, and H.~Lakkaraju,
  ``The disagreement problem in explainable machine learning: A practitioner's perspective,''
  \emph{Transactions on Machine Learning Research (TMLR)}, 2022.

\bibitem{adebayo2018sanity}
J.~Adebayo, J.~Gilmer, M.~Muelly, I.~Goodfellow, M.~Hardt, and B.~Kim, ``Sanity
  checks for saliency maps,'' in \emph{Advances in Neural Information Processing Systems (NeurIPS)}, vol.~31, 2018, pp. 9505--9515.

\bibitem{guo2017calibration}
C.~Guo, G.~Pleiss, Y.~Sun, and K.~Q. Weinberger, ``On calibration of modern
  neural networks,'' in \emph{Proc. International Conference on Machine Learning (ICML)}, 2017, pp. 1321--1330.

\bibitem{hendrycks2017baseline}
D.~Hendrycks and K.~Gimpel, ``A baseline for detecting misclassified and
  out-of-distribution examples in neural networks,'' in \emph{Proc. International Conference on Learning Representations (ICLR)}, 2017.

\bibitem{jocher2023yolo}
G.~Jocher, A.~Chaurasia, and J.~Qiu, ``Ultralytics YOLOv8,'' 2023. [Online]. Available: \url{https://github.com/ultralytics/ultralytics}

\bibitem{kuznetsova2020openimages}
A.~Kuznetsova \emph{et~al.},
  ``The Open Images Dataset V4: Unified image classification, object detection, and visual relationship detection at scale,''
  \emph{International Journal of Computer Vision}, vol. 128, no.~7, pp. 1956--1981, 2020.

\bibitem{dosovitskiy2021vit}
A.~Dosovitskiy \emph{et~al.}, ``An image is worth 16x16 words: Transformers for image recognition at scale,''
  in \emph{Proc. International Conference on Learning Representations (ICLR)}, 2021.

\bibitem{deng2009imagenet}
J.~Deng, W.~Dong, R.~Socher, L.-J. Li, K.~Li, and L.~Fei-Fei, ``ImageNet: A large-scale hierarchical image database,''
  in \emph{Proc. IEEE Conference on Computer Vision and Pattern Recognition (CVPR)}, 2009, pp. 248--255.

\bibitem{lundberg2017unified}
S.~M. Lundberg and S.-I. Lee, ``A unified approach to interpreting model predictions,''
  in \emph{Advances in Neural Information Processing Systems (NeurIPS)}, vol.~30, 2017.

\bibitem{kokhlikyan2020captum}
N.~Kokhlikyan \emph{et~al.}, ``Captum: A unified and generic model interpretability library for PyTorch,''
  \emph{arXiv preprint arXiv:2009.07896}, 2020.

\bibitem{hedstrom2023quantus}
A.~Hedstr{\"o}m \emph{et~al.}, ``Quantus: An explainable AI toolkit for responsible evaluation of neural network explanations,''
  \emph{Journal of Machine Learning Research}, vol.~24, no. 130, pp. 1--11, 2023.

\bibitem{yang2022omnixai}
W.~Yang, H.~Le, S.~C.~H. Hoi, and D.~Sahoo, ``OmniXAI: A library for explainable AI,''
  \emph{Journal of Machine Learning Research}, vol.~23, no. 340, pp. 1--7, 2022.

\bibitem{liu2023visual}
H.~Liu, C.~Li, Q.~Wu, and Y.~J. Lee, ``Visual instruction tuning,''
  in \emph{Advances in Neural Information Processing Systems (NeurIPS)}, vol.~36, 2023.

\bibitem{alayrac2022flamingo}
J.-B. Alayrac \emph{et~al.}, ``Flamingo: A visual language model for few-shot learning,''
  in \emph{Advances in Neural Information Processing Systems (NeurIPS)}, vol.~35, 2022, pp. 23\,716--23\,736.

\bibitem{geminiteam2024gemini}
{Gemini Team, Google}, ``Gemini: A family of highly capable multimodal models,''
  \emph{arXiv preprint arXiv:2312.11805}, 2024.

\bibitem{he2016deep}
K.~He, X.~Zhang, S.~Ren, and J.~Sun, ``Deep residual learning for image recognition,''
  in \emph{Proc. IEEE Conference on Computer Vision and Pattern Recognition (CVPR)}, 2016, pp. 770--778.

\bibitem{achanta2012slic}
R.~Achanta, A.~Shaji, K.~Smith, A.~Lucchi, P.~Fua, and S.~S{\"u}sstrunk,
  ``SLIC superpixels compared to state-of-the-art superpixel methods,''
  \emph{IEEE Transactions on Pattern Analysis and Machine Intelligence}, vol.~34, no.~11, pp. 2274--2282, 2012.

\bibitem{petsiuk2018rise}
V.~Petsiuk, A.~Das, and K.~Saenko, ``RISE: Randomized input sampling for explanation of black-box models,''
  in \emph{Proc. British Machine Vision Conference (BMVC)}, 2018.

\end{thebibliography}

\end{document}
